\documentclass[aspectratio=169,10pt]{beamer}
\usetheme{Madrid}
\usecolortheme{default}
\usepackage[utf8]{inputenc}
\usepackage[T1]{fontenc}
\usepackage{lmodern}
\usepackage{amsmath,amssymb}
\usepackage{enumitem}
\usepackage{xcolor}

\setbeamertemplate{navigation symbols}{}
\setbeamertemplate{footline}[frame number]
\setbeamertemplate{itemize items}[circle]
\setbeamerfont{frametitle}{size=\large,series=\bfseries}

\title{CIA Part 1 -- Questions 414--472}
\subtitle{CAE--Board Relations, Independence Threats, QAIP, Conformance \& External Assessment}
\author{CIA Study Prep}
\date{\today}

\begin{document}
	
	\begin{frame}
		\titlepage
	\end{frame}
	
	\begin{frame}{Instructions}
		\begin{itemize}
			\item 59 questions in two domains only: (A) CAE--board relations and independence threats; (B) QAIP, conformance judgment, and external assessment.
			\item Each question slide is followed by a separate answer/explanation slide.
			\item Every option is defensible; the correct answer turns on a precise Standards priority or a single word.
			\item Read every qualifier (MOST, LEAST, FIRST, BEST, PRIMARY) before answering.
		\end{itemize}
	\end{frame}
	
	% ============================================================
	% DOMAIN A: CAE--BOARD RELATIONS + INDEPENDENCE THREATS (Q414-Q443)
	% ============================================================
	
	% ============================================================
	% Q414
	% ============================================================
	\begin{frame}[t]{Q414}
		\small
		\textbf{Question:} A CAE reports functionally to the audit committee and administratively to the COO. The COO is the subject of an anonymous hotline allegation of misusing corporate aircraft for personal travel. The COO has no authority over the audit plan but controls the CAE's travel budget and staff allocations. The audit committee has not been informed. Which is the MOST appropriate FIRST action?
		\begin{enumerate}[label=\Alph*.,itemsep=1pt,topsep=2pt]
			\item Quietly investigate the allegation personally because the CAE has the strongest fact-finding skills.
			\item Document the independence threat, inform the audit committee of the allegation and the COO's administrative control, and recommend an independent investigator.
			\item Ask the COO to self-report the aircraft usage before escalating.
			\item Refer the matter to the external auditor and take no further action.
		\end{enumerate}
	\end{frame}
	
	\begin{frame}[t]{Q414 -- Answer \& Explanation}
		\small
		\textbf{Correct Answer: B} \\[4pt]
		\textbf{Why:} The COO's administrative control over budget and staff impairs the CAE's objectivity for this specific investigation. Standards 1110/1130/1210.A2 require documenting the threat and escalating to the audit committee with a recommendation for an independent investigator. \\[4pt]
		\textbf{A wrong:} The CAE investigating an administrative superior creates a self-review/conflict threat. \\[2pt]
		\textbf{C wrong:} Asking the subject to self-report tips off the investigation and compromises evidence. \\[2pt]
		\textbf{D wrong:} External audit referral does not discharge internal audit's responsibility.
	\end{frame}
	
	% ============================================================
	% Q415
	% ============================================================
	\begin{frame}[t]{Q415}
		\small
		\textbf{Question:} The audit committee chair asks the CAE to attend all board meetings but not to speak unless the CEO invites a comment. The charter gives the CAE direct and unrestricted access to the board. Which is the MOST accurate statement?
		\begin{enumerate}[label=\Alph*.,itemsep=1pt,topsep=2pt]
			\item The arrangement is acceptable because the chair is the CAE's functional reporting line.
			\item The arrangement restricts the CAE's direct access to the board and impairs independence (Standards 1111/2600); the CAE should request an explicit right to address the board on audit matters.
			\item The arrangement is acceptable if the CEO approves the arrangement in writing.
			\item The arrangement is acceptable because the CAE only reports findings verbally.
		\end{enumerate}
	\end{frame}
	
	\begin{frame}[t]{Q415 -- Answer \& Explanation}
		\small
		\textbf{Correct Answer: B} \\[4pt]
		\textbf{Why:} Direct, unrestricted communication with the board is a core element of independence (Standards 1111/2600). Requiring CEO invitation to speak restricts the CAE's access and impairs independence. \\[4pt]
		\textbf{A wrong:} The chair cannot narrow the CAE's charter-granted access. \\[2pt]
		\textbf{C wrong:} CEO approval does not cure the impairment. \\[2pt]
		\textbf{D wrong:} Verbal reporting does not address the structural restriction.
	\end{frame}
	
	% ============================================================
	% Q416
	% ============================================================
	\begin{frame}[t]{Q416}
		\small
		\textbf{Question:} The audit committee chair asks to meet with the CAE privately, without the CEO. The CEO tells the CAE that such meetings ``undermine management's relationship with the board'' and asks the CAE to decline. Which is the MOST appropriate action?
		\begin{enumerate}[label=\Alph*.,itemsep=1pt,topsep=2pt]
			\item Decline the meeting to preserve the CEO's relationship with the board.
			\item Attend the private meeting; the audit committee chair has authority and private access is a key part of the CAE--board relationship (Standards 1111/2600).
			\item Attend but secretly record the meeting for the CEO.
			\item Decline and report the CEO's request to the external auditor.
		\end{enumerate}
	\end{frame}
	
	\begin{frame}[t]{Q416 -- Answer \& Explanation}
		\small
		\textbf{Correct Answer: B} \\[4pt]
		\textbf{Why:} Private access to the audit committee chair is essential for independence (Standards 1111/2600). A CEO instruction to decline such a meeting is itself a threat to escalate. \\[4pt]
		\textbf{A wrong:} Ceding to the CEO undermines the CAE's independence. \\[2pt]
		\textbf{C wrong:} Secretly recording the meeting is deceptive and improper. \\[2pt]
		\textbf{D wrong:} Reporting to external audit is not the appropriate channel; the audit committee has authority.
	\end{frame}
	
	% ============================================================
	% Q417
	% ============================================================
	\begin{frame}[t]{Q417}
		\small
		\textbf{Question:} The CAE is asked to serve as a voting member of the organization's investment committee, which approves asset allocation decisions. The CAE would receive the same materials and vote as other members. Which is the MOST appropriate action?
		\begin{enumerate}[label=\Alph*.,itemsep=1pt,topsep=2pt]
			\item Accept, because internal audit should contribute to all significant decisions.
			\item Accept only if the CAE does not later audit investment decisions, and disclose the role to the audit committee with safeguards.
			\item Decline the voting role; voting on decisions that internal audit will later audit creates a self-review/management-responsibility threat (Standards 1130/1130.A1).
			\item Accept and document the role in the annual audit report.
		\end{enumerate}
	\end{frame}
	
	\begin{frame}[t]{Q417 -- Answer \& Explanation}
		\small
		\textbf{Correct Answer: C} \\[4pt]
		\textbf{Why:} Voting membership on an investment committee is a management responsibility and creates a self-review threat when internal audit audits investment decisions (Standards 1130/1130.A1). \\[4pt]
		\textbf{A wrong:} Contribution does not override objectivity. \\[2pt]
		\textbf{B wrong:} A safeguard of non-auditing is not sufficient to cure the management-responsibility impairment; advisory membership is preferable. \\[2pt]
		\textbf{D wrong:} Documentation does not cure the impairment.
	\end{frame}
	
	% ============================================================
	% Q418
	% ============================================================
	\begin{frame}[t]{Q418}
		\small
		\textbf{Question:} The CAE was the head of compliance two years ago. The audit committee asks the CAE to lead an audit of the compliance function, including the AML program the CAE implemented. Which is the MOST appropriate action?
		\begin{enumerate}[label=\Alph*.,itemsep=1pt,topsep=2pt]
			\item Accept because two years have passed and the CAE now reports functionally to the audit committee.
			\item Decline to lead the engagement and disclose the self-review threat; assign the engagement to an auditor without prior responsibility, with independent oversight.
			\item Accept but use a smaller sample.
			\item Accept and document the prior role in the workpapers.
		\end{enumerate}
	\end{frame}
	
	\begin{frame}[t]{Q418 -- Answer \& Explanation}
		\small
		\textbf{Correct Answer: B} \\[4pt]
		\textbf{Why:} Standard 1130.A1 prohibits auditing work the auditor managed within the prior year; beyond that, auditing a program the CAE implemented creates a self-review threat requiring reassignment. \\[4pt]
		\textbf{A wrong:} Time passage and reporting lines do not cure self-review. \\[2pt]
		\textbf{C wrong:} Smaller samples do not cure the threat. \\[2pt]
		\textbf{D wrong:} Documentation does not cure the threat.
	\end{frame}
	
	% ============================================================
	% Q419
	% ============================================================
	\begin{frame}[t]{Q419}
		\small
		\textbf{Question:} A new CEO tells the CAE that all audit reports must be delivered to the CEO's office and that the CEO will ``decide what to forward'' to the audit committee. The charter grants the CAE direct access to the audit committee. Which is the MOST appropriate action?
		\begin{enumerate}[label=\Alph*.,itemsep=1pt,topsep=2pt]
			\item Comply, because the CEO is the top executive.
			\item Decline to cede direct access; explain the Standards 1111/2600 requirement and escalate the matter to the audit committee.
			\item Comply but send a copy in parallel so the committee has the report.
			\item Resign.
		\end{enumerate}
	\end{frame}
	
	\begin{frame}[t]{Q419 -- Answer \& Explanation}
		\small
		\textbf{Correct Answer: B} \\[4pt]
		\textbf{Why:} CEO gatekeeping of audit reports impairs the CAE's direct access to the audit committee and violates Standards 1111/2600. Decline and escalate. \\[4pt]
		\textbf{A wrong:} Executive position does not authorize gatekeeping. \\[2pt]
		\textbf{C wrong:} Parallel copies can be filtered/modified; direct access is required. \\[2pt]
		\textbf{D wrong:} Resignation is a last resort.
	\end{frame}
	
	% ============================================================
	% Q420
	% ============================================================
	\begin{frame}[t]{Q420}
		\small
		\textbf{Question:} The audit committee asks the CAE to investigate a fraud allegation against the CEO's son, a senior VP. The CEO controls the CAE's budget. The audit committee chair assures the CAE that ``the committee will back you.'' Which is the MOST appropriate action?
		\begin{enumerate}[label=\Alph*.,itemsep=1pt,topsep=2pt]
			\item Decline because the CEO controls the budget.
			\item Accept, document the threat, ensure the audit committee receives results directly, and consider an independent investigator if the CEO threatens to interfere.
			\item Accept but report only to the chair verbally.
			\item Accept and exclude the son from scope.
		\end{enumerate}
	\end{frame}
	
	\begin{frame}[t]{Q420 -- Answer \& Explanation}
		\small
		\textbf{Correct Answer: B} \\[4pt]
		\textbf{Why:} The audit committee provides functional oversight; safeguards (direct reporting, documentation, independent investigator) mitigate the threat. Standards 1110/1130 support accepting with safeguards. \\[4pt]
		\textbf{A wrong:} Declining leaves a governance-requested engagement undone. \\[2pt]
		\textbf{C wrong:} Verbal-only reporting is not an appropriate channel. \\[2pt]
		\textbf{D wrong:} Excluding the son defeats the engagement.
	\end{frame}
	
	% ============================================================
	% Q421
	% ============================================================
	\begin{frame}[t]{Q421}
		\small
		\textbf{Question:} The CAE is offered an unpaid seat on the board of a charity that receives significant donations from the organization. The charity's board makes decisions on grant recipients, some of which are also customers of the organization. Which is the MOST appropriate action?
		\begin{enumerate}[label=\Alph*.,itemsep=1pt,topsep=2pt]
			\item Accept because the seat is unpaid.
			\item Disclose the opportunity to the audit committee and evaluate whether the relationships create a conflict of interest affecting audit objectivity.
			\item Decline all outside board seats.
			\item Accept and recuse from any audit of the charity.
		\end{enumerate}
	\end{frame}
	
	\begin{frame}[t]{Q421 -- Answer \& Explanation}
		\small
		\textbf{Correct Answer: B} \\[4pt]
		\textbf{Why:} Even unpaid roles can create conflicts when the charity interacts with the organization. Disclosure and evaluation with the audit committee are required (Code Objectivity; Standards 1130). \\[4pt]
		\textbf{A wrong:} Unpaid status does not eliminate the conflict. \\[2pt]
		\textbf{C wrong:} A blanket prohibition is too restrictive. \\[2pt]
		\textbf{D wrong:} Recusal alone may not cure the impairment.
	\end{frame}
	
	% ============================================================
	% Q422
	% ============================================================
	\begin{frame}[t]{Q422}
		\small
		\textbf{Question:} The audit committee has only two members: a former CFO of the organization (retired 18 months ago) and a long-time friend of the CEO. The committee approves the CAE's charter, plan, and compensation. Which is the MOST accurate statement?
		\begin{enumerate}[label=\Alph*.,itemsep=1pt,topsep=2pt]
			\item The committee is independent because its members are not current employees.
			\item The committee's independence is impaired; the former CFO and the CEO's friend create significant threats, and the CAE should document the concern and escalate to the full board (Standards 1110/1111/1130).
			\item Independence is only a concern for external audit.
			\item The CAE should resign.
		\end{enumerate}
	\end{frame}
	
	\begin{frame}[t]{Q422 -- Answer \& Explanation}
		\small
		\textbf{Correct Answer: B} \\[4pt]
		\textbf{Why:} Audit committee independence is impaired when members have recent executive ties or close personal relationships with the CEO. The CAE must document and escalate to the full board (Standards 1110/1111/1130). \\[4pt]
		\textbf{A wrong:} Non-employee status does not equal independence. \\[2pt]
		\textbf{C wrong:} Independence applies to internal audit governance as well. \\[2pt]
		\textbf{D wrong:} Resignation is a last resort; escalation comes first.
	\end{frame}
	
	% ============================================================
	% Q423
	% ============================================================
	\begin{frame}[t]{Q423}
		\small
		\textbf{Question:} The CAE is asked to design and implement a new SOX 404 compliance framework, including selecting controls and writing the control documentation. The CEO argues that internal audit has the strongest control expertise. Which is the MOST appropriate response?
		\begin{enumerate}[label=\Alph*.,itemsep=1pt,topsep=2pt]
			\item Accept because internal audit has the strongest control expertise.
			\item Decline to design and implement the framework; provide advisory observations only, ensuring management owns the design and documentation, and avoid auditing the framework until objectivity is no longer impaired.
			\item Accept and later audit the framework as assurance.
			\item Decline any involvement.
		\end{enumerate}
	\end{frame}
	
	\begin{frame}[t]{Q423 -- Answer \& Explanation}
		\small
		\textbf{Correct Answer: B} \\[4pt]
		\textbf{Why:} Designing and implementing controls is a management responsibility that impairs objectivity and creates self-review (Standards 1130/1130.A1). Internal audit may advise but must not own the framework. \\[4pt]
		\textbf{A wrong:} Expertise does not justify assuming management responsibility. \\[2pt]
		\textbf{C wrong:} Self-review threat: implementing then auditing. \\[2pt]
		\textbf{D wrong:} Advisory involvement is permitted.
	\end{frame}
	
	% ============================================================
	% Q424
	% ============================================================
	\begin{frame}[t]{Q424}
		\small
		\textbf{Question:} The CFO asks the CAE to sign management's SOX 404 assertion that internal control over financial reporting is effective. The CAE has performed an audit and believes controls are effective with two minor documentation gaps. Which is the MOST appropriate action?
		\begin{enumerate}[label=\Alph*.,itemsep=1pt,topsep=2pt]
			\item Sign the assertion because the CAE believes controls are effective.
			\item Decline to sign management's assertion; signing it would assume management responsibility and impair independence. The CAE should report the audit conclusion separately (Standards 1130/1130.A1).
			\item Sign with a note about the documentation gaps.
			\item Sign only if the CEO also signs.
		\end{enumerate}
	\end{frame}
	
	\begin{frame}[t]{Q424 -- Answer \& Explanation}
		\small
		\textbf{Correct Answer: B} \\[4pt]
		\textbf{Why:} Signing management's assertion assumes a management responsibility that impairs independence (Standards 1130/1130.A1). The CAE should report separately. \\[4pt]
		\textbf{A wrong:} Belief does not justify assuming management responsibility. \\[2pt]
		\textbf{C wrong:} A note does not cure the impairment. \\[2pt]
		\textbf{D wrong:} Additional signatories do not cure the impairment.
	\end{frame}
	
	% ============================================================
	% Q425
	% ============================================================
	\begin{frame}[t]{Q425}
		\small
		\textbf{Question:} The CEO offers the CAE's spouse a senior manager position in a business unit that internal audit regularly audits. The CAE learns of the offer before the spouse accepts. Which is the MOST appropriate action?
		\begin{enumerate}[label=\Alph*.,itemsep=1pt,topsep=2pt]
			\item Do nothing; the spouse's employment is a personal matter.
			\item Disclose the offer to the audit committee and evaluate the independence threat; if the spouse accepts, reassign the affected engagements or apply safeguards approved by the full committee.
			\item Ask the CEO to withdraw the offer.
			\item Wait until the spouse accepts, then reassign the engagements.
		\end{enumerate}
	\end{frame}
	
	\begin{frame}[t]{Q425 -- Answer \& Explanation}
		\small
		\textbf{Correct Answer: B} \\[4pt]
		\textbf{Why:} A close family member's employment in an audited area creates a familiarity/self-interest threat (Standard 1130). Disclosure and evaluation are required; reassignment or safeguards follow if the spouse accepts. \\[4pt]
		\textbf{A wrong:} The situation is not purely personal; it affects audit independence. \\[2pt]
		\textbf{C wrong:} The CAE should not interfere with the CEO's hiring decisions. \\[2pt]
		\textbf{D wrong:} Disclosure should occur before acceptance.
	\end{frame}
	
	% ============================================================
	% Q426
	% ============================================================
	\begin{frame}[t]{Q426}
		\small
		\textbf{Question:} The audit committee asks the CAE to audit the audit committee's own processes, including its minutes, charter compliance, and self-assessment. Which is the MOST appropriate action?
		\begin{enumerate}[label=\Alph*.,itemsep=1pt,topsep=2pt]
			\item Decline because the CAE reports to the audit committee.
			\item Accept with safeguards: disclose the self-review threat, use an independent reviewer or external facilitator for the assessment, and ensure objective methodology.
			\item Accept and perform it alone.
			\item Report only verbally.
		\end{enumerate}
	\end{frame}
	
	\begin{frame}[t]{Q426 -- Answer \& Explanation}
		\small
		\textbf{Correct Answer: B} \\[4pt]
		\textbf{Why:} Auditing the audit committee that oversees the CAE creates a self-review/independence threat. Safeguards (independent reviewer, disclosure, objective methodology) are required; the engagement is not prohibited (Standards 1130/1111). \\[4pt]
		\textbf{A wrong:} Declining is unnecessary; safeguards exist. \\[2pt]
		\textbf{C wrong:} Performing it alone without safeguards is improper. \\[2pt]
		\textbf{D wrong:} Verbal-only reporting is insufficient.
	\end{frame}
	
	% ============================================================
	% Q427
	% ============================================================
	\begin{frame}[t]{Q427}
		\small
		\textbf{Question:} The CEO instructs the CAE to cancel a planned audit of a business unit owned by the CEO's family, arguing that the unit is ``not part of the organization.'' The unit uses the organization's shared services (HR, IT, finance) and represents 10\% of consolidated revenue. Which is the MOST appropriate action?
		\begin{enumerate}[label=\Alph*.,itemsep=1pt,topsep=2pt]
			\item Cancel the audit because the CEO has authority.
			\item Keep the audit in the plan based on risk, document the CEO's instruction, and escalate to the audit committee.
			\item Refer the matter to the external auditor and cancel the audit.
			\item Delay the audit until the family situation changes.
		\end{enumerate}
	\end{frame}
	
	\begin{frame}[t]{Q427 -- Answer \& Explanation}
		\small
		\textbf{Correct Answer: B} \\[4pt]
		\textbf{Why:} The unit uses shared services and contributes 10\% of consolidated revenue, creating a material audit universe. A CEO instruction to cancel is an independence threat that must be documented and escalated (Standards 1110/1130/2010/2600). \\[4pt]
		\textbf{A wrong:} CEO authority does not override risk-based coverage and independence. \\[2pt]
		\textbf{C wrong:} External referral does not discharge internal audit's responsibility. \\[2pt]
		\textbf{D wrong:} Delay rewards intimidation and leaves material risk uncovered.
	\end{frame}
	
	% ============================================================
	% Q428
	% ============================================================
	\begin{frame}[t]{Q428}
		\small
		\textbf{Question:} The CAE is asked to consult on M\&A due diligence for a target company. A year later, the CAE is asked to audit the integration of the acquired company. Which is the MOST appropriate action?
		\begin{enumerate}[label=\Alph*.,itemsep=1pt,topsep=2pt]
			\item Accept because a year has passed.
			\item Decline to lead the audit and disclose the self-review threat; assign the audit to an auditor without prior responsibility, with independent oversight.
			\item Accept but use a smaller sample.
			\item Accept and document the prior consultation.
		\end{enumerate}
	\end{frame}
	
	\begin{frame}[t]{Q428 -- Answer \& Explanation}
		\small
		\textbf{Correct Answer: B} \\[4pt]
		\textbf{Why:} Consulting on the acquisition creates a self-review threat for auditing the integration (Standards 1130/1130.A1). Reassignment with independent oversight is required. \\[4pt]
		\textbf{A wrong:} One year does not automatically cure self-review. \\[2pt]
		\textbf{C wrong:} Sample size does not cure the threat. \\[2pt]
		\textbf{D wrong:} Documentation does not cure the threat.
	\end{frame}
	
	% ============================================================
	% Q429
	% ============================================================
	\begin{frame}[t]{Q429}
		\small
		\textbf{Question:} The audit committee chair asks the CAE to investigate a potential conflict of interest involving another board member, but asks the CAE to keep the investigation confidential from the full board until the chair decides how to proceed. Which is the MOST appropriate action?
		\begin{enumerate}[label=\Alph*.,itemsep=1pt,topsep=2pt]
			\item Accept because the chair is the CAE's functional supervisor.
			\item Decline the confidentiality restriction; a board member conflict of interest requires full board awareness and proper governance, not a single-member gatekeeping arrangement.
			\item Accept but document the restriction.
			\item Report the matter directly to the external auditor.
		\end{enumerate}
	\end{frame}
	
	\begin{frame}[t]{Q429 -- Answer \& Explanation}
		\small
		\textbf{Correct Answer: B} \\[4pt]
		\textbf{Why:} A single board member gatekeeping an investigation involving another board member creates a suppression risk. Proper governance requires full board awareness (Standards 1111/1130/2600). \\[4pt]
		\textbf{A wrong:} Functional supervision does not authorize gatekeeping. \\[2pt]
		\textbf{C wrong:} Documentation does not cure the governance failure. \\[2pt]
		\textbf{D wrong:} External auditor is not the appropriate channel.
	\end{frame}
	
	% ============================================================
	% Q430
	% ============================================================
	\begin{frame}[t]{Q430}
		\small
		\textbf{Question:} The CEO asks the CAE to temporarily (6 months) serve as the organization's compliance officer in addition to the CAE role, due to a sudden resignation. The compliance officer reports to the CEO and manages regulatory relationships. Which is the MOST appropriate response?
		\begin{enumerate}[label=\Alph*.,itemsep=1pt,topsep=2pt]
			\item Accept because the arrangement is temporary.
			\item Decline the dual role; serving as compliance officer is a management responsibility and would impair the CAE's independence and objectivity, including for later audits of compliance (Standards 1130/1130.A1).
			\item Accept but recuse from auditing compliance.
			\item Accept and document the dual role.
		\end{enumerate}
	\end{frame}
	
	\begin{frame}[t]{Q430 -- Answer \& Explanation}
		\small
		\textbf{Correct Answer: B} \\[4pt]
		\textbf{Why:} Serving as compliance officer is a management responsibility. Even temporarily, it impairs independence and creates a self-review threat for later audits (Standards 1130/1130.A1). \\[4pt]
		\textbf{A wrong:} Temporary assignments still create impairment. \\[2pt]
		\textbf{C wrong:} Recusal does not cure the impairment for other compliance-related work. \\[2pt]
		\textbf{D wrong:} Documentation does not cure the impairment.
	\end{frame}
	
	% ============================================================
	% Q431
	% ============================================================
	\begin{frame}[t]{Q431}
		\small
		\textbf{Question:} The CFO asks the CAE to approve the CFO's own travel expense report, arguing that the CAE has ``unbiased judgment.'' Which is the MOST appropriate action?
		\begin{enumerate}[label=\Alph*.,itemsep=1pt,topsep=2pt]
			\item Accept because the CAE has unbiased judgment.
			\item Decline; approving the CFO's expenses is a management responsibility and creates a self-review/independence threat. Recommend that approval be handled by another executive or the audit committee.
			\item Accept and document the approval.
			\item Accept only if the amounts are small.
		\end{enumerate}
	\end{frame}
	
	\begin{frame}[t]{Q431 -- Answer \& Explanation}
		\small
		\textbf{Correct Answer: B} \\[4pt]
		\textbf{Why:} Approval of executive expenses is a management function; internal audit performing it creates a management-responsibility/self-review threat (Standards 1130/1130.A1). Recommend an alternative approver. \\[4pt]
		\textbf{A wrong:} Unbiased judgment is not the criterion. \\[2pt]
		\textbf{C wrong:} Documentation does not cure the impairment. \\[2pt]
		\textbf{D wrong:} Amount size is irrelevant to the impairment.
	\end{frame}
	
	% ============================================================
	% Q432
	% ============================================================
	\begin{frame}[t]{Q432}
		\small
		\textbf{Question:} The audit committee asks the CAE to investigate an allegation of CEO misconduct. The CEO is the CAE's administrative reporting line. The audit committee chair has a personal friendship with the CEO. Which is the MOST appropriate action?
		\begin{enumerate}[label=\Alph*.,itemsep=1pt,topsep=2pt]
			\item Accept and report only to the chair.
			\item Accept with safeguards: report to the full audit committee (not the chair alone), document independence considerations, and engage an independent investigator if the chair's friendship impairs oversight.
			\item Decline because the CEO is the administrative reporting line.
			\item Refer to external auditor and close the matter.
		\end{enumerate}
	\end{frame}
	
	\begin{frame}[t]{Q432 -- Answer \& Explanation}
		\small
		\textbf{Correct Answer: B} \\[4pt]
		\textbf{Why:} The chair's friendship with the CEO impairs oversight, so reporting to the full committee and considering an independent investigator is required. Accepting with safeguards is appropriate (Standards 1110/1130/2600). \\[4pt]
		\textbf{A wrong:} Reporting only to the conflicted chair is improper. \\[2pt]
		\textbf{C wrong:} Declining leaves a governance matter unaddressed. \\[2pt]
		\textbf{D wrong:} External referral does not discharge internal audit's responsibility.
	\end{frame}
	
	% ============================================================
	% Q433
	% ============================================================
	\begin{frame}[t]{Q433}
		\small
		\textbf{Question:} The CFO asks the CAE to review the external auditor's fees for ``reasonableness'' before the audit committee approves the fee. Which is the MOST appropriate response?
		\begin{enumerate}[label=\Alph*.,itemsep=1pt,topsep=2pt]
			\item Accept because the CAE is well-positioned to assess fees.
			\item Accept only if the audit committee asks the CAE to do so; the audit committee, not management, has authority over the external auditor relationship and fees.
			\item Decline and report the request to the external auditor.
			\item Accept and provide a specific fee reduction target.
		\end{enumerate}
	\end{frame}
	
	\begin{frame}[t]{Q433 -- Answer \& Explanation}
		\small
		\textbf{Correct Answer: B} \\[4pt]
		\textbf{Why:} The audit committee is responsible for the external auditor relationship and fees. Internal audit can assist only when directed by the audit committee, and should not act on a CFO instruction that bypasses the committee. \\[4pt]
		\textbf{A wrong:} Internal audit's position does not authorize action on a CFO request. \\[2pt]
		\textbf{C wrong:} Reporting to external audit is not the appropriate channel. \\[2pt]
		\textbf{D wrong:} Setting fee reduction targets is a management/committee role.
	\end{frame}
	
	% ============================================================
	% Q434
	% ============================================================
	\begin{frame}[t]{Q434}
		\small
		\textbf{Question:} The CAE was a partner at the organization's external audit firm two years ago. The CAE now leads internal audit and reports functionally to the audit committee. Which is the MOST accurate statement?
		\begin{enumerate}[label=\Alph*.,itemsep=1pt,topsep=2pt]
			\item The prior external audit firm relationship is not relevant to internal audit independence.
			\item The prior relationship may create familiarity and self-interest threats; the CAE should disclose it to the audit committee and ensure safeguards (e.g., no direct coordination with former partners on sensitive judgments) are in place.
			\item The CAE should resign from internal audit.
			\item The prior relationship is irrelevant once the CAE changed employers.
		\end{enumerate}
	\end{frame}
	
	\begin{frame}[t]{Q434 -- Answer \& Explanation}
		\small
		\textbf{Correct Answer: B} \\[4pt]
		\textbf{Why:} Prior employment at the external audit firm can create familiarity/self-interest threats, especially in coordination with the external auditor. Disclosure and safeguards are required (Code Objectivity; Standards 1130). \\[4pt]
		\textbf{A wrong:} The prior relationship is relevant. \\[2pt]
		\textbf{C wrong:} Resignation is an overreaction. \\[2pt]
		\textbf{D wrong:} Changing employers does not eliminate the threat.
	\end{frame}
	
	% ============================================================
	% Q435
	% ============================================================
	\begin{frame}[t]{Q435}
		\small
		\textbf{Question:} The audit committee chair also sits on the compensation committee and personally sets the CAE's bonus. Which is the MOST accurate statement?
		\begin{enumerate}[label=\Alph*.,itemsep=1pt,topsep=2pt]
			\item The arrangement is appropriate because the audit committee chair is independent.
			\item The arrangement may impair objectivity if the CAE's compensation depends on pleasing the same individual who chairs the committee; functional independence is better served when compensation decisions are made by the full committee, not one member.
			\item The arrangement is required by the Standards.
			\item The arrangement is acceptable if the CAE is a CIA.
		\end{enumerate}
	\end{frame}
	
	\begin{frame}[t]{Q435 -- Answer \& Explanation}
		\small
		\textbf{Correct Answer: B} \\[4pt]
		\textbf{Why:} Standard 1111 places CAE performance evaluation and compensation with the audit committee as a body; a single member controlling compensation creates a self-interest/familiarity threat that can compromise objectivity. \\[4pt]
		\textbf{A wrong:} Independence of the chair does not cure the concentration of power. \\[2pt]
		\textbf{C wrong:} The Standards require full-committee oversight, not single-member control. \\[2pt]
		\textbf{D wrong:} Certification does not cure the structural impairment.
	\end{frame}
	
	% ============================================================
	% Q436
	% ============================================================
	\begin{frame}[t]{Q436}
		\small
		\textbf{Question:} The CAE reports administratively to the CFO. The CFO asks the CAE to chair the finance risk committee, which approves finance risk responses. Which is the MOST appropriate response?
		\begin{enumerate}[label=\Alph*.,itemsep=1pt,topsep=2pt]
			\item Accept because the CFO asked and the CAE has risk expertise.
			\item Decline the chair role; chairing a committee that approves risk responses creates a self-review threat when internal audit later assesses finance risks. An advisory/non-voting role may be acceptable with disclosure.
			\item Accept and document the role.
			\item Accept only if the audit committee approves.
		\end{enumerate}
	\end{frame}
	
	\begin{frame}[t]{Q436 -- Answer \& Explanation}
		\small
		\textbf{Correct Answer: B} \\[4pt]
		\textbf{Why:} Chairing a committee that approves risk responses is a management responsibility and creates self-review (Standards 1130/1130.A1). Advisory membership may be acceptable with disclosure. \\[4pt]
		\textbf{A wrong:} Expertise does not justify impaired objectivity. \\[2pt]
		\textbf{C wrong:} Documentation does not cure the impairment. \\[2pt]
		\textbf{D wrong:} Committee approval does not eliminate the self-review threat.
	\end{frame}
	
	% ============================================================
	% Q437
	% ============================================================
	\begin{frame}[t]{Q437}
		\small
		\textbf{Question:} The CAE consulted on the design of a new treasury system eight months ago. The audit committee now asks the CAE to perform an assurance engagement on the treasury system. Which is the MOST appropriate action?
		\begin{enumerate}[label=\Alph*.,itemsep=1pt,topsep=2pt]
			\item Accept because eight months have passed.
			\item Decline to lead the assurance engagement and disclose the self-review threat; assign the engagement to an auditor without prior design responsibility, with independent oversight.
			\item Accept but use a smaller sample.
			\item Accept and document the prior consultation.
		\end{enumerate}
	\end{frame}
	
	\begin{frame}[t]{Q437 -- Answer \& Explanation}
		\small
		\textbf{Correct Answer: B} \\[4pt]
		\textbf{Why:} Consulting on the design creates a self-review threat for assurance (Standards 1130/1130.A1). Reassignment with independent oversight is required. \\[4pt]
		\textbf{A wrong:} Eight months does not automatically cure self-review. \\[2pt]
		\textbf{C wrong:} Sample size does not cure the threat. \\[2pt]
		\textbf{D wrong:} Documentation does not cure the threat.
	\end{frame}
	
	% ============================================================
	% Q438
	% ============================================================
	\begin{frame}[t]{Q438}
		\small
		\textbf{Question:} The board asks the CAE to review the formulation of corporate strategy. The audit committee has already approved the audit plan, which does not include strategy. The CEO objects, saying strategy is ``off-limits'' to internal audit. Which is the MOST appropriate action?
		\begin{enumerate}[label=\Alph*.,itemsep=1pt,topsep=2pt]
			\item Decline because the plan does not include strategy.
			\item Escalate to the audit committee: the board has requested the engagement, so the committee should approve a plan revision; the CEO's objection does not override board direction, but the committee must formally adjust the plan.
			\item Proceed without informing the audit committee.
			\item Refer the matter to the external auditor.
		\end{enumerate}
	\end{frame}
	
	\begin{frame}[t]{Q438 -- Answer \& Explanation}
		\small
		\textbf{Correct Answer: B} \\[4pt]
		\textbf{Why:} Standard 2010 requires the plan to be approved by the board/audit committee. A board request outside the plan requires committee approval and communication. CEO objection does not override board direction. \\[4pt]
		\textbf{A wrong:} Declining ignores a board request. \\[2pt]
		\textbf{C wrong:} Proceeding without committee approval violates Standard 2010. \\[2pt]
		\textbf{D wrong:} External referral does not discharge internal audit's responsibility.
	\end{frame}
	
	% ============================================================
	% Q439
	% ============================================================
	\begin{frame}[t]{Q439}
		\small
		\textbf{Question:} The audit committee chair asks the CAE to attend all executive sessions of the audit committee. The CEO asks the CAE to provide a written summary of each executive session to the CEO. Which is the MOST appropriate action?
		\begin{enumerate}[label=\Alph*.,itemsep=1pt,topsep=2pt]
			\item Provide the summary because the CEO is the administrative reporting line.
			\item Decline to provide the summary; executive sessions are confidential to the committee, and disclosing them to the CEO impairs the CAE's confidential relationship with the committee (Standards 1111/2600).
			\item Provide only a verbal summary.
			\item Provide the summary only if the chair approves.
		\end{enumerate}
	\end{frame}
	
	\begin{frame}[t]{Q439 -- Answer \& Explanation}
		\small
		\textbf{Correct Answer: B} \\[4pt]
		\textbf{Why:} Executive sessions are confidential to the audit committee. Sharing with the CEO would impair the CAE's confidential relationship with the committee (Standards 1111/2600). \\[4pt]
		\textbf{A wrong:} Administrative reporting does not authorize disclosure. \\[2pt]
		\textbf{C wrong:} Verbal sharing is still a breach. \\[2pt]
		\textbf{D wrong:} Even with chair approval, disclosing to management undermines confidentiality.
	\end{frame}
	
	% ============================================================
	% Q440
	% ============================================================
	\begin{frame}[t]{Q440}
		\small
		\textbf{Question:} The CEO tells the CAE that all future communications with the audit committee chair must be pre-approved by the CEO. Which is the MOST appropriate action?
		\begin{enumerate}[label=\Alph*.,itemsep=1pt,topsep=2pt]
			\item Comply because the CEO is the administrative reporting line.
			\item Decline to accept the restriction; direct access to the audit committee chair is essential to independence (Standards 1111/2600), and the CAE should escalate the matter to the audit committee.
			\item Comply but keep a private log of communications.
			\item Resign.
		\end{enumerate}
	\end{frame}
	
	\begin{frame}[t]{Q440 -- Answer \& Explanation}
		\small
		\textbf{Correct Answer: B} \\[4pt]
		\textbf{Why:} Direct access to the audit committee chair is a core independence element (Standards 1111/2600). Pre-approval by the CEO is an unacceptable restriction. \\[4pt]
		\textbf{A wrong:} Administrative reporting does not authorize the restriction. \\[2pt]
		\textbf{C wrong:} A log does not cure the impairment. \\[2pt]
		\textbf{D wrong:} Resignation is a last resort.
	\end{frame}
	
	% ============================================================
	% Q441
	% ============================================================
	\begin{frame}[t]{Q441}
		\small
		\textbf{Question:} The organization's CFO resigns unexpectedly, and the CEO asks the CAE to serve as interim CFO for three months while a search is conducted. Which is the MOST appropriate response?
		\begin{enumerate}[label=\Alph*.,itemsep=1pt,topsep=2pt]
			\item Accept because the interim role is short and the CAE has financial expertise.
			\item Decline the interim CFO role; assuming management responsibility impairs independence and creates a self-review threat for audits of finance. Recommend an alternative interim arrangement and document the decision (Standards 1130/1130.A1).
			\item Accept but recuse from auditing finance.
			\item Accept only if the audit committee approves.
		\end{enumerate}
	\end{frame}
	
	\begin{frame}[t]{Q441 -- Answer \& Explanation}
		\small
		\textbf{Correct Answer: B} \\[4pt]
		\textbf{Why:} Interim CFO is a management responsibility. Even short-term, it impairs independence and creates self-review for future audits (Standards 1130/1130.A1). \\[4pt]
		\textbf{A wrong:} Short duration does not eliminate the impairment. \\[2pt]
		\textbf{C wrong:} Recusal does not cure the impairment of other finance-related audits. \\[2pt]
		\textbf{D wrong:} Committee approval does not cure the impairment.
	\end{frame}
	
	% ============================================================
	% Q442
	% ============================================================
	\begin{frame}[t]{Q442}
		\small
		\textbf{Question:} Internal audit reports administratively to the CEO; the board is passive and meets only twice a year. The CAE wants to strengthen board relations. Which is the MOST appropriate FIRST step?
		\begin{enumerate}[label=\Alph*.,itemsep=1pt,topsep=2pt]
			\item Request a charter revision to establish functional reporting to the audit committee, with the CAE's appointment, evaluation, and compensation subject to committee approval (Standard 1110/1111).
			\item Send monthly reports to the board without audit committee involvement.
			\item Ask the CEO to arrange more board meetings.
			\item Report only to the CEO until the board takes interest.
		\end{enumerate}
	\end{frame}
	
	\begin{frame}[t]{Q442 -- Answer \& Explanation}
		\small
		\textbf{Correct Answer: A} \\[4pt]
		\textbf{Why:} Functional reporting to the audit committee (including appointment, evaluation, and compensation) is required for independence (Standards 1110/1111). A charter revision is the appropriate first step. \\[4pt]
		\textbf{B wrong:} Reports should flow through governance, not directly to the board bypassing the committee. \\[2pt]
		\textbf{C wrong:} Meeting frequency alone does not fix the governance structure. \\[2pt]
		\textbf{D wrong:} Deferring perpetuates the independence weakness.
	\end{frame}
	
	% ============================================================
	% Q443
	% ============================================================
	\begin{frame}[t]{Q443}
		\small
		\textbf{Question:} The audit committee chair asks the CAE to establish a direct feedback channel to senior audit staff for concerns about the CAE's leadership. The CAE supports the idea. Which is the MOST accurate statement?
		\begin{enumerate}[label=\Alph*.,itemsep=1pt,topsep=2pt]
			\item The arrangement is unnecessary because the CAE supports it.
			\item The arrangement strengthens governance and independence by giving the committee independent visibility into the function; it should be documented in the charter or committee protocols.
			\item The arrangement is inappropriate because it undermines the CAE's authority.
			\item The arrangement is optional and purely informal.
		\end{enumerate}
	\end{frame}
	
	\begin{frame}[t]{Q443 -- Answer \& Explanation}
		\small
		\textbf{Correct Answer: B} \\[4pt]
		\textbf{Why:} Committee access to senior audit staff strengthens independence and governance (Standards 1110/1111). Formalizing the channel in committee protocols is the best practice. \\[4pt]
		\textbf{A wrong:} CAE support does not make the channel unnecessary. \\[2pt]
		\textbf{C wrong:} The channel does not undermine the CAE's authority. \\[2pt]
		\textbf{D wrong:} Formalization is more effective than an informal channel.
	\end{frame}
	
	% ============================================================
	% DOMAIN B: QAIP + CONFORMANCE JUDGMENT + EXTERNAL ASSESSMENT (Q444-Q472)
	% ============================================================
	
	% ============================================================
	% Q444
	% ============================================================
	\begin{frame}[t]{Q444}
		\small
		\textbf{Question:} The internal audit activity has a QAIP, but it is not documented. Oversight activities are performed informally, no periodic internal assessments are completed, and the external assessment is 4 years old. Which statement is MOST accurate?
		\begin{enumerate}[label=\Alph*.,itemsep=1pt,topsep=2pt]
			\item The activity conforms because it performs oversight informally.
			\item The activity does not conform; the QAIP must be documented (Standard 1300/1320) and must include periodic internal assessments (Standard 1311).
			\item The external assessment cures all gaps.
			\item Documentation is not required.
		\end{enumerate}
	\end{frame}
	
	\begin{frame}[t]{Q444 -- Answer \& Explanation}
		\small
		\textbf{Correct Answer: B} \\[4pt]
		\textbf{Why:} Standards 1300/1320 require a documented QAIP; Standard 1311 requires periodic internal assessments. Informal oversight alone does not conform. \\[4pt]
		\textbf{A wrong:} Informal oversight is insufficient. \\[2pt]
		\textbf{C wrong:} External assessment does not cure internal QAIP gaps. \\[2pt]
		\textbf{D wrong:} Documentation is required.
	\end{frame}
	
	% ============================================================
	% Q445
	% ============================================================
	\begin{frame}[t]{Q445}
		\small
		\textbf{Question:} The external quality assessment is overdue by 8 months. The CAE argues that internal assessments are thorough and no external assessment is necessary. Which is the MOST accurate statement?
		\begin{enumerate}[label=\Alph*.,itemsep=1pt,topsep=2pt]
			\item The CAE's argument is acceptable if internal assessments are thorough.
			\item The activity does not conform; Standard 1312 requires external assessment at least every five years, and the CAE must disclose the nonconformance and remediate (Standard 1320).
			\item External assessment is only required for large audit functions.
			\item The 8-month overrun is immaterial.
		\end{enumerate}
	\end{frame}
	
	\begin{frame}[t]{Q445 -- Answer \& Explanation}
		\small
		\textbf{Correct Answer: B} \\[4pt]
		\textbf{Why:} Standard 1312 requires external assessment at least every five years. An 8-month overrun is a nonconformance requiring disclosure and remediation (Standard 1320). \\[4pt]
		\textbf{A wrong:} Internal assessments do not substitute for external assessment. \\[2pt]
		\textbf{C wrong:} The requirement applies to all internal audit activities. \\[2pt]
		\textbf{D wrong:} The overrun is a nonconformance regardless of duration.
	\end{frame}
	
	% ============================================================
	% Q446
	% ============================================================
	\begin{frame}[t]{Q446}
		\small
		\textbf{Question:} The external quality assessor's team includes a former internal audit staff member who left the organization six months ago. Which is the MOST accurate statement?
		\begin{enumerate}[label=\Alph*.,itemsep=1pt,topsep=2pt]
			\item The assessment is valid because the team is qualified.
			\item The assessment's independence is impaired; a former staff member who left six months ago creates a familiarity/self-review threat, and the CAE should evaluate and likely obtain a different assessor.
			\item Former staff members are always prohibited.
			\item The organization's independence from the assessor cures the issue.
		\end{enumerate}
	\end{frame}
	
	\begin{frame}[t]{Q446 -- Answer \& Explanation}
		\small
		\textbf{Correct Answer: B} \\[4pt]
		\textbf{Why:} Standard 1312 requires the external assessor to be independent and objective. A former staff member who left six months ago creates a significant independence threat. \\[4pt]
		\textbf{A wrong:} Team qualifications do not cure the independence threat. \\[2pt]
		\textbf{C wrong:} Not always prohibited; the threat must be evaluated. \\[2pt]
		\textbf{D wrong:} Organizational independence does not cure individual threats.
	\end{frame}
	
	% ============================================================
	% Q447
	% ============================================================
	\begin{frame}[t]{Q447}
		\small
		\textbf{Question:} The external quality assessor's firm provides unrelated consulting services to the organization (not to internal audit). Which is the MOST accurate statement?
		\begin{enumerate}[label=\Alph*.,itemsep=1pt,topsep=2pt]
			\item The relationship automatically disqualifies the assessor.
			\item The CAE must evaluate whether the firm's broader relationship with the organization impairs its objectivity for the assessment; unrelated services may not disqualify the firm, but the evaluation must be documented.
			\item The firm is always independent as long as it is not auditing its own work.
			\item Committee approval automatically cures any concern.
		\end{enumerate}
	\end{frame}
	
	\begin{frame}[t]{Q447 -- Answer \& Explanation}
		\small
		\textbf{Correct Answer: B} \\[4pt]
		\textbf{Why:} Standard 1312 requires evaluating objectivity. Unrelated services do not automatically disqualify, but the CAE must document the evaluation of the firm's broader relationship with the organization. \\[4pt]
		\textbf{A wrong:} Not automatically disqualified. \\[2pt]
		\textbf{C wrong:} Independence is not automatic. \\[2pt]
		\textbf{D wrong:} Committee approval does not cure an unaddressed threat.
	\end{frame}
	
	% ============================================================
	% Q448
	% ============================================================
	\begin{frame}[t]{Q448}
		\small
		\textbf{Question:} The CAE concludes the internal audit activity ``conforms'' with the Standards but has not reported the QAIP results or the conformance conclusion to the board. Which is the MOST accurate statement?
		\begin{enumerate}[label=\Alph*.,itemsep=1pt,topsep=2pt]
			\item The conclusion is valid because the CAE performed the QAIP.
			\item The activity does not fully conform; Standards 1320/2430 require reporting QAIP results and conformance to the board at least annually.
			\item Reporting to the CEO is sufficient.
			\item QAIP results are internal only.
		\end{enumerate}
	\end{frame}
	
	\begin{frame}[t]{Q448 -- Answer \& Explanation}
		\small
		\textbf{Correct Answer: B} \\[4pt]
		\textbf{Why:} Standards 1320/2430 require the CAE to report QAIP results and conformance with the Standards to the board at least annually. Omitting it is a nonconformance. \\[4pt]
		\textbf{A wrong:} Performance without reporting is insufficient. \\[2pt]
		\textbf{C wrong:} CEO reporting does not satisfy board reporting. \\[2pt]
		\textbf{D wrong:} QAIP results must be reported to the board.
	\end{frame}
	
	% ============================================================
	% Q449
	% ============================================================
	\begin{frame}[t]{Q449}
		\small
		\textbf{Question:} The QAIP covers engagement performance but not the activity's governance, administration, or resource management. Which is the MOST accurate statement?
		\begin{enumerate}[label=\Alph*.,itemsep=1pt,topsep=2pt]
			\item The QAIP is adequate because engagements are the core of internal audit.
			\item The QAIP is incomplete; Standard 1300 requires the QAIP to cover all aspects of the internal audit activity.
			\item Governance and administration are outside the QAIP.
			\item Resource management is covered by the budget.
		\end{enumerate}
	\end{frame}
	
	\begin{frame}[t]{Q449 -- Answer \& Explanation}
		\small
		\textbf{Correct Answer: B} \\[4pt]
		\textbf{Why:} Standard 1300 requires the QAIP to cover all aspects of the internal audit activity, including governance, administration, and resource management. \\[4pt]
		\textbf{A wrong:} Engagement-only coverage is insufficient. \\[2pt]
		\textbf{C wrong:} Governance and administration are within the QAIP. \\[2pt]
		\textbf{D wrong:} The budget does not substitute for QAIP coverage.
	\end{frame}
	
	% ============================================================
	% Q450
	% ============================================================
	\begin{frame}[t]{Q450}
		\small
		\textbf{Question:} The QAIP's ongoing internal assessments consist only of an annual survey of auditees. No independent post-engagement review is performed. Which is the MOST accurate statement?
		\begin{enumerate}[label=\Alph*.,itemsep=1pt,topsep=2pt]
			\item Auditee surveys alone are a recognized ongoing assessment method.
			\item The ongoing assessments are insufficient; Standard 1311 also expects independent review of engagement work by someone not on the engagement team.
			\item Post-engagement reviews are optional.
			\item Auditee surveys replace post-engagement reviews.
		\end{enumerate}
	\end{frame}
	
	\begin{frame}[t]{Q450 -- Answer \& Explanation}
		\small
		\textbf{Correct Answer: B} \\[4pt]
		\textbf{Why:} Standard 1311 ongoing assessments should include independent review of engagement work (e.g., post-engagement reviews). Surveys alone are insufficient. \\[4pt]
		\textbf{A wrong:} Surveys alone are not sufficient. \\[2pt]
		\textbf{C wrong:} Post-engagement reviews are an expected element. \\[2pt]
		\textbf{D wrong:} Surveys do not replace independent review.
	\end{frame}
	
	% ============================================================
	% Q451
	% ============================================================
	\begin{frame}[t]{Q451}
		\small
		\textbf{Question:} The periodic internal assessment was performed as a self-assessment by the CAE using a recognized maturity model, with no independent validation. Which is the MOST accurate statement?
		\begin{enumerate}[label=\Alph*.,itemsep=1pt,topsep=2pt]
			\item The self-assessment is sufficient because a recognized model was used.
			\item The self-assessment should be validated or supplemented by independent review to support conclusions; self-assessment alone may lack objectivity (Standard 1311).
			\item Periodic internal assessments are optional.
			\item The external assessment alone is sufficient.
		\end{enumerate}
	\end{frame}
	
	\begin{frame}[t]{Q451 -- Answer \& Explanation}
		\small
		\textbf{Correct Answer: B} \\[4pt]
		\textbf{Why:} Standard 1311 requires periodic internal assessments; a recognized model helps, but independent validation enhances objectivity. Self-assessment alone may lack sufficient objectivity. \\[4pt]
		\textbf{A wrong:} A model alone does not ensure objectivity. \\[2pt]
		\textbf{C wrong:} Periodic internal assessments are required. \\[2pt]
		\textbf{D wrong:} External assessment does not replace periodic internal assessments.
	\end{frame}
	
	% ============================================================
	% Q452
	% ============================================================
	\begin{frame}[t]{Q452}
		\small
		\textbf{Question:} The CAE uses QAIP results to determine individual auditor bonuses. Which is the MOST accurate statement?
		\begin{enumerate}[label=\Alph*.,itemsep=1pt,topsep=2pt]
			\item Using QAIP results for bonuses strengthens the QAIP.
			\item Using QAIP results for individual bonuses creates a self-review/independence concern and may reduce the candor of QAIP assessments; the CAE should consider whether this use is appropriate and whether safeguards are needed (Standard 1300/1310).
			\item QAIP results cannot be used for any HR purpose.
			\item The audit committee must approve all QAIP-related bonus decisions.
		\end{enumerate}
	\end{frame}
	
	\begin{frame}[t]{Q452 -- Answer \& Explanation}
		\small
		\textbf{Correct Answer: B} \\[4pt]
		\textbf{Why:} Linking QAIP results to bonuses can create self-review and may compromise the candor of assessments. The CAE should evaluate appropriateness and safeguards (Standards 1300/1310). \\[4pt]
		\textbf{A wrong:} The practice may undermine QAIP integrity. \\[2pt]
		\textbf{C wrong:} QAIP results can inform development and performance management with safeguards. \\[2pt]
		\textbf{D wrong:} Committee approval is not required for all bonus decisions.
	\end{frame}
	
	% ============================================================
	% Q453
	% ============================================================
	\begin{frame}[t]{Q453}
		\small
		\textbf{Question:} One engagement audited a process the CAE designed within the prior year. The CAE concluded the activity ``generally conforms.'' Which is the MOST accurate statement?
		\begin{enumerate}[label=\Alph*.,itemsep=1pt,topsep=2pt]
			\item The conclusion is valid because the engagement was completed.
			\item The activity does not conform; Standard 1130.A1 prohibits auditing work the auditor managed within the prior year, and the nonconformance must be disclosed and remediated (Standard 1320).
			\item The engagement is acceptable because it was documented.
			\item The conclusion is valid because no significant issues were found.
		\end{enumerate}
	\end{frame}
	
	\begin{frame}[t]{Q453 -- Answer \& Explanation}
		\small
		\textbf{Correct Answer: B} \\[4pt]
		\textbf{Why:} Standard 1130.A1 prohibits auditing work the auditor managed within the prior year. The nonconformance must be disclosed and remediated (Standard 1320). \\[4pt]
		\textbf{A wrong:} Completion does not cure the impairment. \\[2pt]
		\textbf{C wrong:} Documentation does not cure the impairment. \\[2pt]
		\textbf{D wrong:} Absence of findings does not cure the impairment.
	\end{frame}
	
	% ============================================================
	% Q454
	% ============================================================
	\begin{frame}[t]{Q454}
		\small
		\textbf{Question:} The external quality assessment team consists of two assessors who both had prior working relationships with the CAE at a previous employer. Which is the MOST accurate statement?
		\begin{enumerate}[label=\Alph*.,itemsep=1pt,topsep=2pt]
			\item The assessment is valid because prior working relationships are common in the profession.
			\item The assessment's independence may be impaired; prior working relationships with the CAE create familiarity threats, and the CAE should evaluate and consider a different assessor team.
			\item The assessment is invalid in all cases.
			\item Committee approval cures any concern.
		\end{enumerate}
	\end{frame}
	
	\begin{frame}[t]{Q454 -- Answer \& Explanation}
		\small
		\textbf{Correct Answer: B} \\[4pt]
		\textbf{Why:} Standard 1312 requires evaluating the external assessor's objectivity. Prior working relationships create familiarity threats that must be evaluated; a different team may be needed. \\[4pt]
		\textbf{A wrong:} Common relationships do not eliminate the threat. \\[2pt]
		\textbf{C wrong:} The assessment is not automatically invalid. \\[2pt]
		\textbf{D wrong:} Committee approval does not cure the threat.
	\end{frame}
	
	% ============================================================
	% Q455
	% ============================================================
	\begin{frame}[t]{Q455}
		\small
		\textbf{Question:} The internal audit activity has no QAIP but claims ``generally conforms'' because engagements follow the Standards. Which is the MOST accurate statement?
		\begin{enumerate}[label=\Alph*.,itemsep=1pt,topsep=2pt]
			\item The claim is valid because engagement-level conformance is the core.
			\item The claim is invalid; Standards 1300-1312 require a QAIP covering all aspects of the activity, and its absence is a systemic nonconformance.
			\item The QAIP is optional for small audit functions.
			\item The external assessment will cure the gap.
		\end{enumerate}
	\end{frame}
	
	\begin{frame}[t]{Q455 -- Answer \& Explanation}
		\small
		\textbf{Correct Answer: B} \\[4pt]
		\textbf{Why:} Standards 1300-1312 require a QAIP for all internal audit activities. Its absence is a systemic nonconformance. \\[4pt]
		\textbf{A wrong:} Engagement conformance does not substitute for the QAIP. \\[2pt]
		\textbf{C wrong:} The QAIP is required for all activities. \\[2pt]
		\textbf{D wrong:} External assessment does not replace the QAIP.
	\end{frame}
	
	% ============================================================
	% Q456
	% ============================================================
	\begin{frame}[t]{Q456}
		\small
		\textbf{Question:} Periodic internal assessments are performed but their results are not reported to senior management or the board. Which is the MOST accurate statement?
		\begin{enumerate}[label=\Alph*.,itemsep=1pt,topsep=2pt]
			\item Reporting is optional because the assessments are internal.
			\item Reporting QAIP results to the board is required (Standards 1320/2430); omitting it is a nonconformance.
			\item The external assessment replaces the need for reporting.
			\item The CAE may report verbally only.
		\end{enumerate}
	\end{frame}
	
	\begin{frame}[t]{Q456 -- Answer \& Explanation}
		\small
		\textbf{Correct Answer: B} \\[4pt]
		\textbf{Why:} Standards 1320/2430 require reporting QAIP results and conformance to the board. Omitting it is a nonconformance. \\[4pt]
		\textbf{A wrong:} Internal-only reporting is insufficient. \\[2pt]
		\textbf{C wrong:} External assessment does not replace CAE reporting. \\[2pt]
		\textbf{D wrong:} Verbal-only reporting does not satisfy the requirement.
	\end{frame}
	
	% ============================================================
	% Q457
	% ============================================================
	\begin{frame}[t]{Q457}
		\small
		\textbf{Question:} The board requests QAIP results. The CAE states that QAIP results are ``internal to the audit function'' and declines to share them. Which is the MOST accurate statement?
		\begin{enumerate}[label=\Alph*.,itemsep=1pt,topsep=2pt]
			\item The CAE's position is correct; QAIP results are internal.
			\item The CAE's position is incorrect; Standards 1320/2430 require reporting QAIP results and conformance to the board.
			\item The board must request results through the CEO.
			\item QAIP results are confidential to the audit committee only.
		\end{enumerate}
	\end{frame}
	
	\begin{frame}[t]{Q457 -- Answer \& Explanation}
		\small
		\textbf{Correct Answer: B} \\[4pt]
		\textbf{Why:} Standards 1320/2430 require QAIP results and conformance reporting to the board. Refusal is a nonconformance. \\[4pt]
		\textbf{A wrong:} QAIP results must be reported to the board. \\[2pt]
		\textbf{C wrong:} The CAE reports directly to the board. \\[2pt]
		\textbf{D wrong:} Reporting is to the board/audit committee, not confidential to one group.
	\end{frame}
	
	% ============================================================
	% Q458
	% ============================================================
	\begin{frame}[t]{Q458}
		\small
		\textbf{Question:} The external quality assessment was completed 5 years and 1 month ago. No new assessment has been started. Which is the MOST accurate statement?
		\begin{enumerate}[label=\Alph*.,itemsep=1pt,topsep=2pt]
			\item The activity conforms because 5 years is roughly satisfied.
			\item The activity does not conform; Standard 1312 requires external assessment at least every five years, and the CAE must disclose the nonconformance and remediate (Standard 1320).
			\item A one-month overrun is immaterial.
			\item The next external assessment can be deferred indefinitely.
		\end{enumerate}
	\end{frame}
	
	\begin{frame}[t]{Q458 -- Answer \& Explanation}
		\small
		\textbf{Correct Answer: B} \\[4pt]
		\textbf{Why:} Standard 1312 requires external assessment at least every five years. Exceeding the five-year window is a nonconformance requiring disclosure and remediation (Standard 1320). \\[4pt]
		\textbf{A wrong:} ``Roughly'' does not satisfy the requirement. \\[2pt]
		\textbf{C wrong:} The overrun is a nonconformance. \\[2pt]
		\textbf{D wrong:} Deferral is not permitted.
	\end{frame}
	
	% ============================================================
	% Q459
	% ============================================================
	\begin{frame}[t]{Q459}
		\small
		\textbf{Question:} The activity has: (1) a board-approved charter; (2) a QAIP with ongoing and periodic internal assessments; (3) an external assessment completed within five years; (4) all engagements documented and supervised; (5) CAE reports functionally to the audit committee. Which is the MOST accurate statement?
		\begin{enumerate}[label=\Alph*.,itemsep=1pt,topsep=2pt]
			\item The activity does not conform because external assessment is required annually.
			\item The activity appears to conform with the Standards; the CAE should still verify that QAIP results are reported to the board and that no other nonconformances exist.
			\item The activity does not conform because periodic internal assessments are optional.
			\item The activity cannot conform without a mature analytics program.
		\end{enumerate}
	\end{frame}
	
	\begin{frame}[t]{Q459 -- Answer \& Explanation}
		\small
		\textbf{Correct Answer: B} \\[4pt]
		\textbf{Why:} The listed elements align with the Standards. The CAE must still verify reporting to the board and check for other nonconformances before concluding ``conforms.'' \\[4pt]
		\textbf{A wrong:} External assessment is required every five years, not annually. \\[2pt]
		\textbf{C wrong:} Periodic internal assessments are required. \\[2pt]
		\textbf{D wrong:} Analytics maturity is not a conformance requirement.
	\end{frame}
	
	% ============================================================
	% Q460
	% ============================================================
	\begin{frame}[t]{Q460}
		\small
		\textbf{Question:} The internal audit charter is approved by the CEO, not the board. The CAE claims conformance because the CEO is ``senior management.'' Which is the MOST accurate statement?
		\begin{enumerate}[label=\Alph*.,itemsep=1pt,topsep=2pt]
			\item The claim is valid because the CEO is senior management.
			\item The claim is invalid; Standard 1000 requires the charter to be approved by the board (or its audit committee); the nonconformance must be disclosed and remediated (Standard 1320).
			\item The claim is valid if the CEO also approves the audit plan.
			\item The charter is not a conformance requirement.
		\end{enumerate}
	\end{frame}
	
	\begin{frame}[t]{Q460 -- Answer \& Explanation}
		\small
		\textbf{Correct Answer: B} \\[4pt]
		\textbf{Why:} Standard 1000 requires board approval of the charter. CEO approval is a nonconformance requiring disclosure and remediation (Standard 1320). \\[4pt]
		\textbf{A wrong:} Senior management approval does not satisfy Standard 1000. \\[2pt]
		\textbf{C wrong:} Plan approval does not substitute for charter approval. \\[2pt]
		\textbf{D wrong:} The charter is a core conformance requirement.
	\end{frame}
	
	% ============================================================
	% Q461
	% ============================================================
	\begin{frame}[t]{Q461}
		\small
		\textbf{Question:} Engagement workpapers were destroyed six months after the report was issued. Organization policy requires 7-year retention; regulatory requirements require 5 years. Which is the MOST accurate statement?
		\begin{enumerate}[label=\Alph*.,itemsep=1pt,topsep=2pt]
			\item The destruction is acceptable because the report was issued.
			\item The activity does not conform; Standards 2330/2330.A1 require retention per organizational policy and legal requirements, and the premature destruction violates both.
			\item The 5-year regulatory requirement governs, so 6 months is acceptable.
			\item The CAE should recreate the workpapers.
		\end{enumerate}
	\end{frame}
	
	\begin{frame}[t]{Q461 -- Answer \& Explanation}
		\small
		\textbf{Correct Answer: B} \\[4pt]
		\textbf{Why:} Standards 2330/2330.A1 require retention per organizational policy and legal requirements. Destroying workpapers six months after issuance violates both the 7-year policy and the 5-year regulatory minimum. \\[4pt]
		\textbf{A wrong:} Report issuance does not authorize destruction. \\[2pt]
		\textbf{C wrong:} Both requirements apply; the more stringent governs. \\[2pt]
		\textbf{D wrong:} Recreation does not cure the violation.
	\end{frame}
	
	% ============================================================
	% Q462
	% ============================================================
	\begin{frame}[t]{Q462}
		\small
		\textbf{Question:} The external quality assessment concluded ``does not conform.'' The CAE reported the result to the CEO but not to the audit committee. Which is the MOST accurate statement?
		\begin{enumerate}[label=\Alph*.,itemsep=1pt,topsep=2pt]
			\item Reporting to the CEO is sufficient.
			\item The CAE violated Standard 1320 by not reporting the significant nonconformance to the board; the CEO report does not satisfy the requirement.
			\item The external assessor is responsible for reporting to the board.
			\item The report can be withheld if the CAE disagrees with the finding.
		\end{enumerate}
	\end{frame}
	
	\begin{frame}[t]{Q462 -- Answer \& Explanation}
		\small
		\textbf{Correct Answer: B} \\[4pt]
		\textbf{Why:} Standard 1320 requires the CAE to report significant nonconformance to senior management and the board. CEO-only reporting is insufficient. \\[4pt]
		\textbf{A wrong:} CEO reporting does not satisfy board reporting. \\[2pt]
		\textbf{C wrong:} The CAE is responsible for reporting. \\[2pt]
		\textbf{D wrong:} Disagreement does not justify withholding.
	\end{frame}
	
	% ============================================================
	% Q463
	% ============================================================
	\begin{frame}[t]{Q463}
		\small
		\textbf{Question:} The QAIP includes recommendations from assessments but no process for monitoring whether they are implemented. Which is the MOST accurate statement?
		\begin{enumerate}[label=\Alph*.,itemsep=1pt,topsep=2pt]
			\item The QAIP is complete because assessments and recommendations exist.
			\item The QAIP is incomplete; it must include monitoring of whether recommendations are implemented and reporting results to the board (Standards 1300/1320).
			\item Monitoring recommendations is optional.
			\item The external assessment covers monitoring.
		\end{enumerate}
	\end{frame}
	
	\begin{frame}[t]{Q463 -- Answer \& Explanation}
		\small
		\textbf{Correct Answer: B} \\[4pt]
		\textbf{Why:} Standard 1300 requires the QAIP to cover all aspects of the activity, including follow-up on recommendations and reporting to the board (Standard 1320). \\[4pt]
		\textbf{A wrong:} Recommendations without follow-up are incomplete. \\[2pt]
		\textbf{C wrong:} Monitoring is required. \\[2pt]
		\textbf{D wrong:} External assessment does not cover ongoing monitoring.
	\end{frame}
	
	% ============================================================
	% Q464
	% ============================================================
	\begin{frame}[t]{Q464}
		\small
		\textbf{Question:} The CAE reports conformance to the board annually but does not disclose a significant internal assessment finding on engagement supervision. Which is the MOST accurate statement?
		\begin{enumerate}[label=\Alph*.,itemsep=1pt,topsep=2pt]
			\item The annual report is sufficient because conformance is reported.
			\item The CAE violated Standard 1320 by not disclosing the significant nonconformance; reporting conformance without disclosing known significant nonconformance is misleading.
			\item Minor findings do not require disclosure.
			\item The external assessment would have caught the issue.
		\end{enumerate}
	\end{frame}
	
	\begin{frame}[t]{Q464 -- Answer \& Explanation}
		\small
		\textbf{Correct Answer: B} \\[4pt]
		\textbf{Why:} Standard 1320 requires reporting significant nonconformance. Reporting ``conforms'' while omitting a known significant nonconformance is misleading. \\[4pt]
		\textbf{A wrong:} Partial reporting is misleading. \\[2pt]
		\textbf{C wrong:} Significant findings require disclosure. \\[2pt]
		\textbf{D wrong:} The external assessment does not excuse non-disclosure.
	\end{frame}
	
	% ============================================================
	% Q465
	% ============================================================
	\begin{frame}[t]{Q465}
		\small
		\textbf{Question:} The charter is board-approved, but the CAE reports administratively to the CFO with no functional reporting line to the audit committee. Which is the MOST accurate statement?
		\begin{enumerate}[label=\Alph*.,itemsep=1pt,topsep=2pt]
			\item The arrangement conforms because the charter is board-approved.
			\item The arrangement does not conform; Standard 1110/1111 requires functional reporting to the board/audit committee (including appointment, evaluation, and compensation), and the absence impairs independence.
			\item Administrative reporting to the CFO is sufficient.
			\item The arrangement is acceptable if the CFO is a CPA.
		\end{enumerate}
	\end{frame}
	
	\begin{frame}[t]{Q465 -- Answer \& Explanation}
		\small
		\textbf{Correct Answer: B} \\[4pt]
		\textbf{Why:} Standards 1110/1111 require functional reporting to the board/audit committee. Absence of functional reporting impairs independence. \\[4pt]
		\textbf{A wrong:} Charter approval does not substitute for functional reporting. \\[2pt]
		\textbf{C wrong:} CFO reporting is administrative, not functional. \\[2pt]
		\textbf{D wrong:} CPA status is irrelevant.
	\end{frame}
	
	% ============================================================
	% Q466
	% ============================================================
	\begin{frame}[t]{Q466}
		\small
		\textbf{Question:} The QAIP's ongoing internal assessments are conducted only at the engagement level, not at the activity level (e.g., governance, administration, resources). Which is the MOST accurate statement?
		\begin{enumerate}[label=\Alph*.,itemsep=1pt,topsep=2pt]
			\item The QAIP is adequate because engagement-level assessments cover the main work.
			\item The QAIP is incomplete; Standard 1300 requires the QAIP to cover all aspects of the internal audit activity, and ongoing assessments should address activity-level matters as well.
			\item Activity-level assessments are covered by the external assessment.
			\item Only engagement-level assessments are required.
		\end{enumerate}
	\end{frame}
	
	\begin{frame}[t]{Q466 -- Answer \& Explanation}
		\small
		\textbf{Correct Answer: B} \\[4pt]
		\textbf{Why:} Standard 1300 requires QAIP coverage of all aspects of the activity. Ongoing assessments should include activity-level matters, not just engagements. \\[4pt]
		\textbf{A wrong:} Engagement-level alone is insufficient. \\[2pt]
		\textbf{C wrong:} External assessment does not cover ongoing activity-level assessments. \\[2pt]
		\textbf{D wrong:} Activity-level assessments are required.
	\end{frame}
	
	% ============================================================
	% Q467
	% ============================================================
	\begin{frame}[t]{Q467}
		\small
		\textbf{Question:} The external quality assessment team includes a former internal audit staff member who was terminated three months ago for performance reasons. Which is the MOST accurate statement?
		\begin{enumerate}[label=\Alph*.,itemsep=1pt,topsep=2pt]
			\item The assessment is valid because the team is qualified.
			\item The former staff member creates a significant independence threat, particularly given the termination; the CAE should evaluate and likely obtain a different assessor team (Standard 1312).
			\item Termination history is irrelevant to independence.
			\item Former employees are always allowed on the assessment team.
		\end{enumerate}
	\end{frame}
	
	\begin{frame}[t]{Q467 -- Answer \& Explanation}
		\small
		\textbf{Correct Answer: B} \\[4pt]
		\textbf{Why:} Standard 1312 requires the external assessor to be independent and objective. A recently terminated staff member creates a significant self-review/familiarity threat. \\[4pt]
		\textbf{A wrong:} Qualifications do not cure the independence threat. \\[2pt]
		\textbf{C wrong:} Termination history is highly relevant. \\[2pt]
		\textbf{D wrong:} Not automatically allowed; the threat must be evaluated.
	\end{frame}
	
	% ============================================================
	% Q468
	% ============================================================
	\begin{frame}[t]{Q468}
		\small
		\textbf{Question:} The external assessment concluded ``generally conforms'' with one ``partially conforms'' finding on engagement supervision. The CAE claims ``conforms.'' Which is the MOST accurate statement?
		\begin{enumerate}[label=\Alph*.,itemsep=1pt,topsep=2pt]
			\item ``Conforms'' is valid because the overall conclusion was ``generally conforms.''
			\item The ``partially conforms'' finding indicates a nonconformance; the CAE must evaluate its impact, disclose if significant, and remediate (Standards 1320/2340), not claim ``conforms.''
			\item ``Partially conforms'' is not a nonconformance.
			\item The external assessor's conclusion is binding.
		\end{enumerate}
	\end{frame}
	
	\begin{frame}[t]{Q468 -- Answer \& Explanation}
		\small
		\textbf{Correct Answer: B} \\[4pt]
		\textbf{Why:} A ``partially conforms'' finding indicates a nonconformance (Standard 2340 supervision). The CAE must evaluate impact, disclose if significant, and remediate (Standard 1320). \\[4pt]
		\textbf{A wrong:} ``Generally conforms'' with a partial nonconformance is not full conformance. \\[2pt]
		\textbf{C wrong:} It is a nonconformance. \\[2pt]
		\textbf{D wrong:} The CAE evaluates the conclusion.
	\end{frame}
	
	% ============================================================
	% Q469
	% ============================================================
	\begin{frame}[t]{Q469}
		\small
		\textbf{Question:} The QAIP results are not reported to the board, but the external assessment report was forwarded to the board. Which is the MOST accurate statement?
		\begin{enumerate}[label=\Alph*.,itemsep=1pt,topsep=2pt]
			\item The activity conforms because the board received the external assessment.
			\item The activity does not fully conform; Standards 1320/2430 require the CAE to report QAIP results and conformance to the board, and forwarding the external report does not satisfy that responsibility.
			\item External reporting satisfies QAIP reporting.
			\item QAIP results are internal.
		\end{enumerate}
	\end{frame}
	
	\begin{frame}[t]{Q469 -- Answer \& Explanation}
		\small
		\textbf{Correct Answer: B} \\[4pt]
		\textbf{Why:} Standards 1320/2430 require the CAE to report QAIP results and conformance to the board. Forwarding the external report does not satisfy the CAE's reporting responsibility. \\[4pt]
		\textbf{A wrong:} External reporting is not a substitute. \\[2pt]
		\textbf{C wrong:} The CAE must report QAIP results directly. \\[2pt]
		\textbf{D wrong:} QAIP results must be reported to the board.
	\end{frame}
	
	% ============================================================
	% Q470
	% ============================================================
	\begin{frame}[t]{Q470}
		\small
		\textbf{Question:} The CAE argues that no QAIP is needed because the internal audit activity is ``small.'' Which is the MOST accurate statement?
		\begin{enumerate}[label=\Alph*.,itemsep=1pt,topsep=2pt]
			\item The argument is valid because small functions cannot support a QAIP.
			\item The argument is invalid; Standards 1300-1312 require a QAIP for all internal audit activities regardless of size, though the formality and scale may be proportionate.
			\item The QAIP is optional for functions with fewer than five auditors.
			\item External assessment is not required for small functions.
		\end{enumerate}
	\end{frame}
	
	\begin{frame}[t]{Q470 -- Answer \& Explanation}
		\small
		\textbf{Correct Answer: B} \\[4pt]
		\textbf{Why:} Standards 1300-1312 require a QAIP for all internal audit activities. Scale and formality may vary, but the requirement does not. \\[4pt]
		\textbf{A wrong:} Size does not waive the requirement. \\[2pt]
		\textbf{C wrong:} There is no size-based exception. \\[2pt]
		\textbf{D wrong:} External assessment is required regardless of size.
	\end{frame}
	
	% ============================================================
	% Q471
	% ============================================================
	\begin{frame}[t]{Q471}
		\small
		\textbf{Question:} The internal audit activity has internal and external assessments but no follow-up on findings and no monitoring of corrective actions. Which is the MOST accurate statement?
		\begin{enumerate}[label=\Alph*.,itemsep=1pt,topsep=2pt]
			\item The QAIP is complete because assessments exist.
			\item The QAIP is incomplete; it must include follow-up on findings and monitoring of corrective actions, with reporting to the board (Standards 1300/1320).
			\item Follow-up is optional.
			\item The external assessment covers follow-up.
		\end{enumerate}
	\end{frame}
	
	\begin{frame}[t]{Q471 -- Answer \& Explanation}
		\small
		\textbf{Correct Answer: B} \\[4pt]
		\textbf{Why:} Standard 1300 requires the QAIP to cover all aspects of the activity, including follow-up on findings and corrective actions; Standard 1320 requires reporting to the board. \\[4pt]
		\textbf{A wrong:} Assessments without follow-up are incomplete. \\[2pt]
		\textbf{C wrong:} Follow-up is required. \\[2pt]
		\textbf{D wrong:} External assessment does not cover ongoing follow-up.
	\end{frame}
	
	% ============================================================
	% Q472
	% ============================================================
	\begin{frame}[t]{Q472}
		\small
		\textbf{Question:} The external assessment is 4 years old and internal assessments are clean. The CAE claims ``conforms.'' Which is the MOST accurate statement?
		\begin{enumerate}[label=\Alph*.,itemsep=1pt,topsep=2pt]
			\item The claim is valid because internal assessments are clean.
			\item The external assessment is within the five-year requirement (Standard 1312), so the age is not itself a nonconformance; the CAE must still verify that all other conformance requirements are met and that QAIP results are reported to the board.
			\item The 4-year-old external assessment is a nonconformance.
			\item The claim is invalid because external assessment is required annually.
		\end{enumerate}
	\end{frame}
	
	\begin{frame}[t]{Q472 -- Answer \& Explanation}
		\small
		\textbf{Correct Answer: B} \\[4pt]
		\textbf{Why:} 4 years is within the five-year window (Standard 1312). Age is not the issue; the CAE must verify all other conformance requirements, including board reporting. \\[4pt]
		\textbf{A wrong:} Internal assessments alone are insufficient, though not the issue here. \\[2pt]
		\textbf{C wrong:} 4 years is within the window. \\[2pt]
		\textbf{D wrong:} External assessment is required every five years, not annually.
	\end{frame}
	
	\begin{frame}
		\centering
		\Large End of Questions 414--472\\[8pt]
		\normalsize Domains: CAE--Board Relations \& Independence Threats; QAIP, Conformance Judgment \& External Assessment.
	\end{frame}
	
\end{document}